\subsection{选题意义}

首先，本课题推进了低时延推理执行模型的认识边界。现有工作多以扩大融合范围、消除 host 侧 launch 为目标；Megakernel 把依赖判断下沉到设备端后，真正决定临界路径的，往往是算子边的就绪粒度、任务投放时机与物理执行粒度是否匹配。将同一逻辑计算图中的相邻边纳入统一分析，可以解释为何有的边应保留细粒度流水，有的边反而适合用物理 grid 完成语义隔离，从而把优化关键推进到“依赖应由哪一级机制实现”，并建立这类问题的理论模型。

其次，本课题为任意计算图上的 Megakernel 构建提供了系统化路径。面向推理场景，仅有局部手工拆边或单一 persistent grid 并不足以覆盖多样依赖形态；需要一套从依赖分析、任务拆解到执行计划生成的完整框架：先形成可调度的任务图，再按边选择同一 grid 内的事件同步、内部任务调度器，或跨 grid 完成语义，并在选定物理边界上配置 PDL 触发与消费者独立前缀。该路径使图级融合、设备端调度与硬件依赖执行落在同一决策空间，便于形成可复现、可验证的运行时方案。

再次，本课题在 Megakernel 框架内把 grid 边界上的 PDL 与依赖分析结合起来，进一步整合调度时机与调度方式。依赖分析不仅回答一条边应保留在同一 grid 内，还是映射为跨 grid 完成语义，还进一步约束后继何时具备调度资格、以何种机制放行：同一 grid 内可由事件或就绪门控控制投放；跨 grid 时，除串行等待生产者完成外，还可通过 PDL 在 programmatic completion 后提前获得调度资格，并先执行合法独立前缀\cite{nvidia2026pdl,nvidia2026cutlassDependentLaunch}。由此，边界选择、触发时机与交接方式进入同一探索空间，使“何时调度”和“怎样调度”成为依赖感知执行计划的一体两面。

最后，本课题具有直接的服务与工程价值。交互式 LLM 推理对单步时延敏感，局部临界路径上数十微秒的差异会沿 token 序列累积为可感知的 TPOT 与响应延迟。将上述分析固化为可嵌入编译器与运行时的完整系统，有望为低 batch decode、短增量 prefill 及智能体快速返回等场景提供可落地的时延收益。
